(1) 半径均为 \(a\) 的两平行刚性圆盘之间夹有不可压Newton黏液，黏度 \(\mu\)，瞬时均匀间距 \(h(t)\ll a\)，下盘固定、上盘以速度 \(\dot h<0\) 纯法向挤压。忽略惯性、边缘弯月面和空化，边缘压强为 \(p_a\)。由轴对称Stokes平衡与不可压连续方程严格推导主导阶径向速度 \(u_r(r,z)\)、径向流率 \(q_r\) 及Reynolds方程。
(2) 积分求 \(p(r)-p_a\) 的分布和支持的法向载荷 \(F\)，确定中心压强及其与 \(h\) 的标度，证明 \(F\) 对匀速挤压始终为正。
(3) 上盘受到固定向下净外力 \(F_0>0\)，质量与惯性忽略。求 \(h(t)\) 的闭式和任意正厚度 \(h_1<h_0\) 所需时间，判断有限时间内能否达到严格零间隙，检查初始 \(t=0\)。
(4) 用径向黏性剪切 \(\mu(\partial_z u_r)^2\) 在两盘间的体积积分独立求总耗散功率，证明其等于载荷乘闭合速度；检验量纲，说明当 \(h\) 逼近分子尺度、表面滑移或空化出现时应如何理解奇异结论。
